%%
%% This is LaTeX2e input.
%%

%% The following tells LaTeX that we are using the 
%% style file amsart.cls (That is the AMS article style
%%
\documentclass{amsart}
\usepackage{fancyvrb}
\usepackage{amsmath}
\usepackage{amsfonts}
\usepackage{amsthm}
\usepackage{cases}
\usepackage{caption}
\usepackage{xcolor}
\usepackage{tikz}
\usepackage{centernot}
\usepackage{mathtools}
\usepackage[margin=1in]{geometry}
\usepackage[makeroom]{cancel}
\usetikzlibrary{decorations.pathreplacing}
\newcommand{\tikznode}[3][inner sep=0pt]{\tikz[remember
picture,baseline=(#2.base)]{\node(#2)[#1]{$#3$};}}

\usepackage{ytableau}
%\usepackage{witharrows}

\newtheorem{lemma}{Lemma}[section]
\newtheorem{corollary}[lemma]{Corollary}
\newtheorem{conjecture}[lemma]{Conjecture} 
\newtheorem{proposition}[lemma]{Proposition} 
\newtheorem{theorem}[lemma]{Theorem}
\newtheorem{question}{Question}
%%
%% If some other type is need, say conjectures, then it is constructed
%% by editing and uncommenting the following.
%%

%\newtheorem{conj}[thm]{Conjecture} 


%%% 
%%% The following gives definition type environments (which only differ
%%% from theorem type invironmants in the choices of fonts).  The
%%% numbering is still tied to the theorem counter.
%%% 

\theoremstyle{definition}
\newtheorem{definition}{Definition}[section]
\newtheorem{example}{Example}[section]

%%
%% Again more of these can be added by uncommenting and editing the
%% following. 
%%

%\newtheorem{note}[thm]{Note}


%%% 
%%% The following gives remark type environments (which only differ
%%% from theorem type invironmants in the choices of fonts).  The
%%% numbering is still tied to the theorem counter.
%%% 


\theoremstyle{remark}

\newtheorem{remark}[lemma]{Remark}


%%%
%%% The following, if uncommented, numbers equations within sections.
%%% 

\numberwithin{equation}{section}


%%%
%%% The following show how to make definition (also called macros or
%%% abbreviations).  For example to use get a bold face R for use to
%%% name the real numbers the command is \mathbf{R}.  To save typing we
%%% can abbreviate as

\newcommand{\R}{\mathbf{R}}  % The real numbers.

%%
%% The comment after the defintion is not required, but if you are
%% working with someone they will likely thank you for explaining your
%% definition.  
%%
%% Now add you own definitions:
%%

%%%
%%% Mathematical operators (things like sin and cos which are used as
%%% functions and have slightly different spacing when typeset than
%%% variables are defined as follows:
%%%

\DeclareMathOperator{\dist}{dist} % The distance.



%%
%% This is the end of the preamble.
%% 
\title{Triple quantum Schubert calculus}
\author{Matthew J. Samuel}
%\ead{matthematics@gmail.com}

%\affiliation[1]{addressline={31 Cherry Blossom Drive},
%city={Monroe Township},
%postcode={NJ 08831},
%country={USA}}


\DeclareMathOperator{\bpath}{BPath}
\DeclareMathOperator{\Sym}{Sym}
\DeclareMathOperator{\Schub}{Schub}
\DeclareMathOperator{\ESchub}{ESchub}
\DeclareMathOperator{\sch}{\mathfrak{S}}
\newcommand{\tom}[1]{\xrightarrow{#1}}
%\newcommand{\tomm}[2]{\xrightarrow{#1}_{#2}}
%\counterwithout{equation}{section}
%\setcounter{section}{-1}

\begin{document}

\begin{abstract}
\end{abstract}
\maketitle

{\bf Keywords}: Schubert polynomial, structure constant, Littlewood-Richardson rule, coproduct, double Schubert polynomial

{\bf Mathematics classification codes}: 05E05, 05A05, 05A15, 14N10

{\bf Email}: matthematics@gmail.com

\newcommand{\code}{\mathfrak{c}}
\newcommand{\ducode}{\overline{\mathfrak{c}}}
%$$w_0(n)^{(i)}=w_0(n)w_0(n+1-i)w_0(n-i)$$

%\begin{theorem}
%Let $u\in S_n$ and $1\leq i\leq n$. Then we have
%$$\sch_u(y)=\sum_{u'w_0(n)^{(i)}\tom{n+1-i} uw_0(n)}y_i^{\ell(u',u)}\sch_{u'}(y^{(i)})$$
%where $u'\in S_n$ satisfies $\ell(u'w_0(n)^{(i)})+\ell(u')=\ell(w_0(n)^{(i)})$.
%\end{theorem}
\section{Introduction}


\section{Definitions and statement of results}

\begin{definition}[$\tom{k}$, $P_k(u,w)$, $q_k(u,w)$, $n_k(u,w)$] \label{definition:pieri}
We define a relation $\tom{k}$ on pairs of permutations, introduced in its non-quantum form by Sottile \cite{sottile}.%, which was used originally to indicate when the coefficient of a Schubert polynomial is nonzero in a suitable application of the Pieri formula. In the generalization to double Schubert polynomials in \cite{samuelmolev}, the same relation applies to all nonzero coefficients in the expansion of the generalized Pieri formula. The relation plays the same role in our case.

Given $u,u'\in S_\infty$ and a positive integer $k$, we declare that $u\tom{k} w$ if there exists a $p$ with $0\leq p\leq k$ such that there are transpositions $t_{a_1b_1},\ldots,t_{a_pb_p}$ satisfying 
\begin{enumerate}
\item $a_i\leq k<b_i$ for all $i$
\item $a_i\neq a_j$ whenever $i\neq j$
\item $b_i\geq b_j$ whenever $i<j$
\item We have either
$$\ell(ut_{a_1b_1}\cdots t_{a_ib_i})=\ell(ut_{a_1b_1}\cdots t_{a_{i-1}b_{i-1}})+1$$ 
or 
$$\ell(ut_{a_1b_1}\cdots t_{a_ib_i})=\ell(ut_{a_1b_1}\cdots t_{a_{i-1}b_{i-1}})-2(b_i-a_i)+1$$
for all $i$
\item $w=ut_{a_1b_1}\cdots t_{a_pb_p}$
\end{enumerate}
If $u\tom{k} w$, define 
$$P_k(u,w)=\{u(i)\mid i\leq k\mbox{ and }u(i)=w(i)\}$$
Also, we define $q_k(u,w)$ recursively. If $u=w$, define $q_k(u,w)=1$. If we have defined $q_k(u,w')$ with $w'=t_{a_1b_1}\cdots t_{a_{p-1}b_{p-1}}$, suppose $\ell(w't_{a_pb_p})=\ell(w')+1$. Then define
$$q_k(u,w)=q_k(u,w')$$
If instead $\ell(w't_{a_pb_p})=\ell(w')-2(b_p-a_p)+1$, define
$$q_k(u,w)=q_k(u,w')q_{a_pb_p}$$
We also define
$$n_k(u,w)=\#\{i\mid i\leq k\mbox{ and }u(i)\neq w(i)\}$$
\end{definition}

\begin{theorem}[The Pieri formula] \label{theorem:pieri}
Suppose $u\in S_\infty$ and $p\leq k$. Then
$$\sch_u^q(x;y)E_{p,k}^q(x;z)=\sum_{u\tom{k} w}q_k(u,w)E_{p-n_k(u,w),k-n_k(u,w)}(y_{P_k(u,w)};z)\sch_w^q(x;y)$$
\end{theorem}

\begin{theorem} \label{theorem:coeffreduce}
Suppose $u,v\in S_\infty$ and let $d$ be a weak composition such that $d_i\in \{0,1\}$ for all $i$ and if $d_i=1$, then $d_{i+1}=0$. Let $i_1,\ldots,i_k$ be the integers such that $d_{i_p}=1$ for $1\leq p\leq k$. Then $c_{u,v}^{w,d}(y;z)=0$ unless $\ell(us_{i_p})<\ell(u)$ and $\ell(vs_{i_p})<\ell(v)$ for each $p$, in which case
$$c_{u,v}^{w,d}(y;z)=c_{us',vs'}^w(y;z)$$
where
$$s'=s_{i_1}\cdots s_{i_k}$$
\end{theorem}

\newcommand{\lom}{\widetilde{\lambda}}
\newcommand{\mom}{\widetilde{\mu}}
\newcommand{\thop}{\widetilde{\theta}}
\newcommand{\Tom}[1]{\xRightarrow{#1}}
\newcommand{\dom}{\mu}

\begin{definition}[$\code$, (strict) dominant permutations, $\theta$, $\dom$, $\thop$, $\mom$]
We say that a permutation $\mu$ is \emph{dominant} if $\code_i(\mu)\geq \code_{i+1}(\mu)$ for all $i$. A permutation $\mu$ is said to be \emph{strict dominant} if $\mu$ is dominant and whenever $\code_i(\mu)\neq 0$ we have $\code_i(\mu)>\code_{i+1}(\mu)$. For a permutation $v\in S_\infty$, we define a dominant permutation $\dom(v)$ recursively as follows. If $v$ is already dominant, define $\mu(v)=v$. Otherwise, let $i$ be the maximal index such that $\code_i(v)<\code_{i+1}(v)$, and define
$$\dom(v)=\dom(vs_i)$$
Define $\theta_i(v)$ to be $\code_i(\dom(v))$.

Similarly, for a permutation $v\in S_\infty$, we define a strict dominant permutation $\mom(v)$ recursively as follows. If $v$ is already strict dominant, define $\mom(v)=v$. Otherwise, let $i$ be the maximal index such that $\code_{i+1}(v)\neq 0$ and $\code_i(v)\leq \code_{i+1}(v)$, then define
$$\mom(v)=\mom(vs_i)$$
Define $\thop_i(v)$ to be $\code_i(\mom(v))$.
\end{definition}

\begin{definition}[$\xRightarrow{i}$, $\sigma_i(u,v)$, $D_i(u,v)$]
Let $u,v\in S_\infty$ and $i\geq 1$ be an integer. We define $u\Tom{i} v$ if there exist distinct integers $b_1,b_2,\ldots,b_k$ such that $b_j\neq i$ for all $j$, $\ell(ut_{i,b_1}\cdots t_{i,b_j})=\ell(u)+j$ for all $j$, and 
$$v=ut_{i,b_1}\cdots t_{i,b_k}$$
We define
$$\sigma_i(u,v)=\#\{j\mid b_j<i\}$$
and we define
$$D_i(u,v)=\{u(i)\}\cup\{u(j)\mid u(j)\neq v(j)\}$$
\end{definition}

\begin{theorem}[Formula for ordinary double Schubert polynomials] \label{theorem:orddouble}
Suppose $v\in S_\infty$. Then
$$\sch_v(x;z)=\sum_{(v_0,\ldots,v_k)}\prod_{i=1}^k (-1)^{\sigma_i(v_{i-1},v_i)}E_{\theta_i(v)-\ell(v_{i-1},v_i),\theta_i(v)}(x;z_{D_i(v_{i-1},v_i)})$$
Where $v_{i-1}\xRightarrow{i} v_i$ for all $i$, $v_0=1$, and $v_k=v\dom(v^{-1})$.
\end{theorem}


\begin{theorem}[Formula for quantum double Schubert polynomials] \label{theorem:quandouble}
Suppose $v\in S_\infty$. Then
$$\sch_v^q(x;z)=\sum_{(v_0,\ldots,v_k)}\prod_{i=1}^k (-1)^{\sigma_i(v_{i-1},v_i)}E_{\thop_i(v)-\ell(v_{i-1},v_i),\thop_i(v)}^q(x;z_{D_i(v_{i-1},v_i)})$$
Where $v_{i-1}\xRightarrow{i} v_i$ for all $i$, $v_0=1$, and $v_k=v\mom(v^{-1})$.
\end{theorem}

Coupled with the Pieri formulas, these formulas give a relatively efficient method for computing products of (quantum) double Schubert polynomials.




\begin{definition}[$\chi_{j;n}$, $\chi^{(n)}$]
Recall the definition of the commuting difference operators $\chi_{j;n}$, with $j\leq n$, defined by Fomin, Gelfand, and Postnikov, as follows:
$$\chi_{j,n}=x_j-\sum_{1\leq i<j}q_{ij}\partial^{t_{ij}}+\sum_{j<k\leq n} q_{jk}\partial^{t_{jk}}$$
We denote the sequence $(\chi_{1;n},\ldots,\chi_{n;n})$ with $n$ fixed by $\chi^{(n)}$. Recall (CITATION) that if $\sch_v(x)$ has at most $n$ variables, then
$$\sch_v^q(\chi^{(n)})(1) = \sch_v(x)$$
and hence
$$\sch_v^q(\chi^{(n)};z)(1) = \sch_v(x;z)$$
We specify a more precise formula, as follows.
\end{definition}

\newcommand{\qsch}{\sch^q}

\begin{theorem}
Suppose $v\in S_n$. Then
$$\sch_v^q(\chi^{(n)};z)=\sum_{u\in S_{n},d} q^dc_{u,v}^{1,d}(x;z)\partial^u$$
\end{theorem}
\begin{proof}
Note that if $\sch_u(x)$ has at most $n$ variables, then
$$\sch_u^q(\chi^{(n)};y)\sch_v^q(\chi^{(n)};z)=\sum_{w,d} q^dc_{u,v}^{w,d}(y;z)\sch_w^q(\chi^{(n)};y)$$
We thus have that
$$(\sch_u^q(\chi^{(n)};y)\sch_v^q(\chi^{(n)};z))(1)=\sum_{w,d} q^dc_{u,v}^{w,d}(y;z)\sch_w(x;y)$$
This is also equal to
$$\qsch_v(\chi^{(n)};z)(\sch_u(x;y))$$
due to the fact that $(\sch_u^q(\chi^{(n)};y))(1)=\sch_u(x;y)$. To obtain the coefficient of $\partial^u$ in $\qsch_v(\chi^{(n)};z)$, we set $y=x$ after applying this operator. Since $\sch_w(x;x)=0$ unless $w=1$, we have that
$$\qsch_v(\chi^{(n)};z)(\sch_u(x;y))|_{y=x}=\sum_d q^dc_{u,v}^{1,d}(x;z)$$
and the result follows.
\end{proof}

%\begin{lemma}
%Let $k$ be an integer and let $v\in S_\infty$. Then
%$$\sum_d q^dc_{s_k,v}^{1,d}(x;z) = \sum_{i\leq k < j} q_{ij}\sch_{vt_{ij}}(x;z)$$
%\end{lemma}
%\begin{proof}
%This value is equal to
%$$\qsch_{s_k}(\chi^{(n)};y)(\sch_v(x;z))$$
%then setting $y=x$ for sufficiently large $n$. We also have
%$$\qsch_{s_k}(\chi^{(n)};y) = \sch_{s_k}(x;y)+\sum_{1\leq i\leq k<j\leq n}q_{ij}\partial^{t_{ij}}$$
%Thus
%$$\qsch_{s_k}(\chi^{(n)};y)(\sch_v(x;z))=\sum_w c_{s_k,v}^w(y;z)\sch_w(x;y)+\sum_{1\leq i\leq k<j\leq n}q_{ij}\sch_{vt_{ij}}(x;z)$$
%Setting $y=z$, the terms involving $\sch_w(x;y)$ disappear, and we are left with the desired result.
%\end{proof}

\begin{theorem}
Let $k$ be an integer and let $v,w\in S_\infty$. Then
$$\sum_d q^dc_{s_k,v}^{w,d}(y;z)=c_{s_k,v}^{w}(y;z)+\sum_{i\leq k < j} q_{ij}\sch_{vt_{ij}w^{-1}}(y;z)$$
\end{theorem}
\begin{proof}
We have that the desired value is equal to the coefficient of $\sch_w(x;y)$ in
$$\qsch_{s_k}(\chi^{(n)};y)(\sch_{v}(x;z))$$
We have that
$$\qsch_{s_k}(\chi^{(n)};y)=\sch_{s_k}(x;y)+\sum_{i\leq k<j\leq n}q_{ij}\partial^{t_{ij}}$$
Applying this operator, we obtain
$$\sum_{i\leq k<j\leq n} q_{ij}\sch_{vt_{ij}}(x;z)+\sum_w c_{s_k,v}^w(y;z)\sch_w(x;y)$$
Applying the Cauchy formula, this is equal to
$$\sum_{i\leq k<j,u} q_{ij}\sch_{vt_{ij}u^{-1}}(y;z)\sch_u(x;y)+\sum_w c_{s_k,v}^w(y;z)\sch_w(x;y)$$
Setting $u=w$ and allowing $n$ to be sufficiently large, we obtain the result.
\end{proof}

\begin{proof}[Proof of Theorem \ref{theorem:coeffreduce}]
It suffices to prove this for $v=w_0(n)$. We have
$$\qsch_{w_0(n)}(\chi;z)=\sum_{u,d} q^dc_{u,w_0(n)}^{1,d}(x;z)\partial^u$$
Fixing the particular $d$ in the statement of the theorem, we must have that $u=s'$. By the Pieri formula,
$$c_{s',w_0(n)}^{1,d}(x;z)=\sch_{w_0(n)s'}(x;z)$$
Thus if we compute the coefficient of $q^d\sch_w(x;y)$ in 
$$(\qsch_{w_0(n)}(\chi;z))(\sch_u(x;y))$$
we obtain
$$c_{us',w_0(n)s'}^w(y;z)$$
and the result follows.
\end{proof}

\section{Proofs}

\begin{lemma}
%Suppose $u,v,w$ are such that $u\Tom{i} w$, $v\Tom{i} w$, and $D_i(u,w)=D_i(v,w)$. Then $u=v$.
Suppose $(v_0,v_1,\ldots,v_n)$ and $(v_0',\ldots,v_n')$ are such that $v_{i-1}\Tom{i} v_i$, $v_{i-1}'\Tom{i} v_i'$ for all $i$, $v_0=v_0'=1$, and $v_n=v_n'=v$ for some $n$. Suppose also that $D_i(v_{i-1},v_i)=D_{i}(v_{i-1}',v_i')$ for all $i$. Then $v_i=v_i'$ for all $i$.
\end{lemma}
\begin{proof}
This is true if $\ell(v_n)=0$. Otherwise, consider the right descent. There are only two places where it can occur.
\end{proof}


%Acting on $321$-avoiding, all the elem syms are Schubs.

\begin{theorem} \label{theorem:nocan}
No two terms in Theorem \ref{theorem:orddouble} directly cancel each other.
\end{theorem}
\begin{proof}
%Let $(v_0,\ldots,v_m)$ and $(v_0',\ldots,v_m')$ with $v_0=v_0'=1$ and $v_m=v_m'=v\mu(v^{-1})$ be paths such that $v_{i-1}\Tom{i} v_i$ and $v_{i-1}'\Tom{i} v_i'$ for all $i$. It will suffice to show that if 
%$$D_i(v_{i-1},v_i)=D_i(v_{i-1}',v_i')$$
%for all $i$, then $v_i = v_i'$ for all $i$. We prove this is true for all $i\geq j$ by induction on $j$. This is clear if $j=m$. Otherwise, assume it holds for $j+1$. Then we know that $v_{j+1}=v_{j+1}'$. Since $D_{j+1}(v_{j},v_{j+1})=D_{j+1}(v_{j}',v_{j+1}')$, we have that $v_j(k)\neq v_{j+1}(k)$ if and only if $v_j'(k)\neq v_{j+1}'(k)$. Let $a_1,\ldots,a_p$ be the set of indices such that $v_j(a_k),v_j'(a_k)\neq v_{j+1}(a_k)$ for all $k$. Assume the order is such that $v_j^{-1}v_{j+1}(a_k)=a_{k+1}$, with $a_0=a_{p+1}=j$.
\end{proof}

\section{Proof of the Pieri formula}

\begin{lemma} \label{lemma:singlevariable}
Let $u\in S_\infty$ and let $i,j\geq 1$. Then
$$(x_i-z_j)\sch_u^q(x;y)=(y_{u(i)}-z_j)\sch_u^q(x;y)+\sum_{\substack{p\neq i\\\ell(ut_{ip})=\ell(u)+1}}\mathrm{sign}(p-i)\sch_{ut_{ip}}^q(x;y)$$
$$+\sum_{\substack{p\neq i\\\ell(ut_{ip})=\ell(u)-2|p-i|+1}}\mathrm{sign}(p-i)q_{ip}\sch_{ut_{ip}}^q(x;y)$$
\end{lemma}
%\begin{proof}
%It can be shown by the recurrence relation for $\partial_u^w$ that $\partial_u^u=u$. Thus we have
%$$\partial_u^u(x_i-z_j)=x_{u(i)}-z_j$$
%Substituting $y$ for $x$, we obtain that the coefficient of $\sch_u(x;y)$ is $y_{u(i)}-z_j$.

%To compute the constant coefficients we use the equivariant Chevalley formula
%$$\sch_u(x;y)\sch_{s_i}(x;y)=(\sch_{s_i}(u(x);y))|_{x=y}\sch_u(x;y)+\sum_{\substack{u\tom{i} u'\\\ell(u')=\ell(u)+1}}\sch_{u'}(x;y)$$
%We have that
%$$x_i-z_j=\sch_{s_i}(x;y)-\sch_{s_{i-1}}(x;y)+y_i-z_j$$
%with the negative term $\sch_{s_{i-1}}(x;y)$ excluded if $i=1$. Thus the terms occurring with constant coefficients in the expansion of $(x_i-z_j)\sch_u(x;y)$ are
%$$\sum_{\substack{u\tom{i} u'\\\ell(u,u')=1}}\sch_{u'}(x;y)-\sum_{\substack{u\tom{i-1}u''\\\ell(u,u'')=1}}\sch_{u''}(x;y)$$
%The overlap in the two sums is when $u'=u''=ut_{qk}$ where $i\notin \{q,k\}$, which cancel, and what remains is the negative terms $ut_{qi}$ where $q<i$ and the positive terms $ut_{iq}$ where $q>i$. Therefore,
%$$\partial_{u}^{ut_{qk}}(x_i-z_j)=0$$
%if $i\notin \{q,k\}$, and if $q\neq i$ then
%$$\partial_u^{ut_{iq}}(x_i-z_j)=\mathrm{sign}(q-i)$$
%The result follows.
%\end{proof}

\newcommand{\ntom}[1]{\centernot{\tom{#1}}}

\newcommand{\desc}{\mathrm{desc}}

\begin{definition}[Descents of a cycle, $\desc_u(c)$]
Let $c$ be a cycle and suppose $u\in S_\infty$. Write $c$ as
$$(a_1,\ldots,a_n)$$
where $a_n>a_i$ for all $i<n$ (this is a unique way to express the cycle since the $a_j$ are distinct). Then for any $i$ with $1\leq i<n$ we say that $a_i$ is a \emph{descent} of $c$ with respect to $u$ if $u(a_i)>u(a_{i+1})$. We denote the set of all descents of $c$ with respect to $u$ by $\desc_u(c)$.
\end{definition}

\begin{lemma} \label{lemma:piericycle}
Suppose $u,w\in S_\infty$ and let $k>0$ be an integer. Then $u\tom{k} w$ if and only if there exist pairwise disjoint cyclic permutations $c_1,\ldots,c_n$ of respective cycle lengths $p_1+1,\ldots,p_n+1$ such that
$$w=uc_1\cdots c_n$$
and
$$\ell(u,w)=p_1+\cdots+p_n$$
such that for each $i$ there exist pairwise distinct positive integers $a_{i,1},\ldots,a_{i,p_i}$ and $b_i$ such that we may write the cycle $c_i$ as
$$c_i=(a_{i,p_i},a_{i,p_i-1},\ldots, a_{i,1},b_i)$$
where $a_{i,j}\leq k$ for all $1\leq j\leq p_i$, $b_i>k$, and $b_i>b_j$ if $i<j$, and for each $m$ we have
$$\ell(uc_1\cdots c_m) = \sum_{i=1}^m \left(p_i - 2\sum_{a\in \desc_u(c_i)}(b_i-a)\right)$$
\end{lemma}
\begin{proof}
Suppose first that $u\tom{k} w$. Then there exist reflections $t_{a_1b_1},\ldots,t_{a_mb_m}$ for some $m$ with $a_i\neq a_j$ if $i\neq j$, $a_i\leq k<b_i$ for all $i$, $\ell(ut_{a_1b_1}\cdots t_{a_ib_i})=\ell(t_{a_1b_1}\cdots t_{a_{i-1}b_{i-1}})+1$ or $\ell(ut_{a_1b_1}\cdots t_{a_ib_i})=\ell(t_{a_1b_1}\cdots t_{a_{i-1}b_{i-1}})-2(b_i-a_i)+1$ for all $i$, $b_i\geq b_j$ if $i\leq j$, and $w=ut_{a_1b_1}\cdots t_{a_mb_m}$. Fix $b$ such that there is an index $1\leq p\leq m$ such that $b=b_p$ and let $i$ and $r$ be such that $a_i,\ldots,a_{i+r}$ is the maximal portion of the sequence such that $b_j=b$ for all $i\leq j\leq i+r$. Let
$$c=t_{a_ib_i}\cdots t_{a_{i+r}b_{i+r}}$$
Given $j$ with $i\leq j\leq i+r$ we claim that if we define a permutation $c^{(j)}$ with
$$c^{(j)}=t_{a_ib}\cdots t_{a_jb}$$
then
$$c^{(j)}=(a_j,a_{j-1},\ldots,a_i,b)$$
and that
$$\ell(uc^{(j)})=\ell(u)+j-2\sum_{a\in\desc_u(c^{(j)})}(b-a)$$
This is clear if $j=i$. Otherwise, suppose
$$c^{(j-1)}=t_{a_ib}\cdots t_{a_{j-1}b}=(a_{j-1},\ldots,a_i,b)$$
Then
$$c^{(j-1)}t_{a_jb}=(a_{j-1},\ldots,a_i,b)(b,a_j)=(a_{j-1},\ldots,a_i,b,a_j)=(a_j,\ldots,a_i,b)$$
as desired. We have either that 
$$u(a_j)=uc^{(j-1)}(a_j)<uc^{(j-1)}(b)=u(a_{j-1})$$
and $\ell(uc^{(j)})=\ell(uc^{(j-1)})+1$, or
$$u(a_j)>u(a_{j-1})$$
and $a_j\in\desc_u(c^{(j)})$ and $\ell(uc^{(j)})=\ell(uc^{(j-1)})-2(a_j-b)+1$. Hence the desired claim follows by induction. Hence,
$$c=c^{(i+r)}=(a_{i+r},\ldots,a_i,b)$$
and the length satisfies the condition. Ranging over all $b$, these cycles, which we index in order of decreasing $b$ as $c_1,\ldots,c_n$, are the ones that occur in the disjoint cycle decomposition of $u^{-1}w$, which we can see by the fact that the cycles are pairwise disjoint and $w=uc_1\ldots c_n$. 

We prove the converse now. Suppose the cycles $c_1,\ldots,c_n$ of cycle length $p_i+1$ for each $i$ exist satisfying the conditions, fixing $k$. Write $c_1$ as
$$c_1=(a_{p_1},\ldots,a_1,b)$$
Then
$$c_1=t_{a_1b}\cdots t_{a_{p_1}b}$$
We claim that $\ell(ut_{a_1b}\cdots t_{a_ib})>\ell(ut_{a_1b}\cdots t_{a_{i-1}b})$ if $a_i\notin\desc_u(c_1)$, and $\ell(ut_{a_1b}\cdots t_{a_ib})<\ell(ut_{a_1b}\cdots t_{a_{i-1}b})$ if $a_i\in\desc_u(c_1)$. This clearly holds for $a_1$. Otherwise, we note that
$$ut_{a_1b}\cdots t_{a_{i-1}b}(b)=u(a_{i-1})$$
and
$$ut_{a_1b}\cdots t_{a_{i-1}b}(a_i)=u(a_{i})$$
since $a_j\neq a_i$ if $j\neq i$. Thus, if $a_i\notin \desc_u(c_1)$, then the length increases, and if $a_i\in \desc_u(c_1)$ then the length decreases. Applying this argument for each disjoint cycle $c_i$ to $\ell(uc_1\cdots c_i)$ and expanding the cycles into reflections, we obtain that successively multiplying on the right by each reflection increases the length by at least $1$ for each $a_i$ that is not a descent and decreases the length by at most $2(b_i-a)-1$ for each descent in $c_i$. By the length condition in the statement of the lemma, the increases must be by exactly $1$ and the decreases by exactly $2(b_i-a)-1$ by maximality.
%Since there are $p_1+\cdots+p_n$ reflections and this is equal to $\ell(u,w)$, each reflection must increase the length by exactly $1$ and we have the result.
\end{proof}

\newcommand{\lup}{\triangleleft}

\begin{definition}[$\lup$]
If $u,u'\in S_\infty$, we define $u\lup u'$ if $\ell(u,u')=1$ or $u'=ut_{ab}$ where $\ell(u,u')=1-2(b-a)$.

%Define $\nabla_u^w(F(x))$ to be the coefficient of $\sch_w^q(x;y)$ in the product $F(x)\sch_u^q(x;y)$, replacing $y$ with $x$ and quantizing.

%Same as unquantizing, doing the div diff, and requantizing.

%$$\mathcal{Q}\partial^w\mathcal{Q}^{-1}$$
%$$\mathcal{Q}\partial^w\mathcal{Q}^{-1}(\sch_u^q(x;y)\sch_v^q(x;z))$$
%
%$$\mathcal{Q}\partial^w\mathcal{Q}^{-1}(c_{u,v}^{w,d}(y;z)q^d\sch_w^q(x;y))$$
%
%$$\mathcal{Q}\partial^wc_{u,v}^{w',d}(y;z)q^d\sch_{w'}(x;y)$$
%
%$$c_{u,v}^{w',d}(y;z)q^d\sch_{w'w^{-1}}^q(x;y)$$
%Product
%$$\nabla^{s_i}p = \sum c_{u,v}^{s_i,d}(x;x)q^d\nabla^u(p)\nabla^v$$
%
%Ideal:
%$$p-\sum_u c_{u,v}^{1,d}(x;x)q^d\nabla^u(p)\nabla^v$$
%$$(\sch^q_u(\chi;y)\sch_v^q(\chi;y))\cdot 1=c_{u,v}^w(x;x)\sch_w(y;x)$$
%$$F(x)\sch_u^q(x;y)=\sum_w \nabla_u^w(F(x))|_{x=y}\sch_w^q(x;y)$$
%$$\nabla^u(\sch_v^q(x;z))=\nabla^u(\sch_{v_2}^q(x;y)\sch_{v_1}(y;z)$$
\end{definition}



\begin{lemma} \label{lemma:onedescent}
Suppose $u,w\in S_\infty$ and let $k>1$ be an integer. Suppose $u\tom{k}w$ and let $c=(a_p,\ldots,a_1,b)$ be a cycle as in Lemma \ref{lemma:piericycle} in the decomposition of $u^{-1}w$. Then $\desc_u(c)$ has at most one element.
\end{lemma}
\begin{proof}
This is clearly true if $p=1$, or if the result holds for $p-1$ but $c^{(p-1)}$ has no descents. Suppose then that $c^{(p-1)}$ has a descent, say $a_i$ for exactly one $i<p$ by the induction hypothesis, and that $u(a_p)>u(a_{p-1})$, meaning that $a_p\in\desc_u(c^{(p-1)})$. $a_i$ being a descent would require that 
$$u(a_i)>uc^{(i-1)}(j)>u(a_{i-1})$$
for all $j$ with $a_i<j<b$. Furthermore, we also have that for all $j$ with $a_p<j<b$ that
$$u(a_p)>uc^{(p-1)}(j)>u(a_{p-1})$$
Suppose $a_p<a_i$. Then $u(a_p)>uc^{(p-1)}(a_i)=u(a_{i-1})>u(a_{p-1})$. This means that necessarily $a_{p-1}<a_i$, since $u(a_{i-1})=uc^{(i)}(a_i)$ is minimal in the range $a_i\leq j\leq b$ in $uc^{(i)}$. For the same reason, we cannot have that $uc^{(p-2)}(b)<u(a_{i-1})$, because this would have to come from some $t_{a_qb}$ with $i<q<p-1$, $a_q<a_i$, and and $u(a_q)<u(a_{i-1})$, which would increase the length by more than $1$. However, since $uc^{(p-2)}(b)>u(a_{i-1})>u(a_{p-1})$, this would imply that $t_{a_{p-1}b}$ increases the length by more than $1$, which is a contradiction. Thus we cannot have a descent $a_p$ with $a_p<a_i$. 

If instead $a_p>a_i$, then $u(a_i)>u(a_p)>u(a_{i-1})$. Since $a_p$ is a descent, $uc^{(p-1)}(b)=u(a_{p-1})<u(a_p)$, but $uc^{(i)}(b)=u(a_i)>u(a_p)$. This would imply that there is some $q$ with $i<q\leq p-1$ such that $a_q>a_p$ and $uc^{(q-1)}t_{a_qb}(b)=u(a_q)<u(a_p)$, because if we had $a_q<a_p$ then the length would increase by more than $1$. Let $q$ be minimal with this property. In that case, $u(a_{q-1})>u(a_p)$ and $uc^{(p-1)}(a_q)=u(a_{q-1})$, contradicting the fact that $a_p$ is a descent. It follows by induction that there can be at most one descent.
\end{proof}

\begin{definition}[Redefinition of $\desc_u(c)$]
If $c$ is a cycle, we redefine $\desc_u(c)$ to be the unique element of the previously defined set $\desc_u(c)$ if it exists, or otherwise $0$.
\end{definition}

\begin{lemma} \label{lemma:geq}
Suppose $u,w\in S_\infty$ and let $k>1$ be an integer. Suppose $u\tom{k}w$ and let $c=(a_p,\ldots,a_1,b)$ be a cycle as in Lemma \ref{lemma:piericycle} in the decomposition of $u^{-1}w$, with $a_q=k$ being a descent. Then for all $j\geq q$ we have that $u(a_j)>u(b)$.
\end{lemma}
\begin{proof}
We have that $u(a_q)>u(a_{q-1})$. By Lemma \ref{lemma:onedescent}, $k$ is the only descent of $c$. We claim that for all $r$ with $k\leq r<b$ we have that $u(r)>u(b)$. This is because for all such $r$ we have that $u(a_{q-1})<u(r)$ since $a_q$ is a descent, and if $u(b)>u(r)$ for some such $r$ then we may pick a maximal $\ell<q$ such that $u(a_{\ell})>u(r)$ for some such $r$. Then $u(a_{\ell+1})<u(r)$ for all such $r$, and the reflection $t_{a_{\ell-1} a_{\ell}}$ would increase the length by more than $1$, which cannot happen. For the same reason, we cannot have that $u(a_m)<u(r)$ for any $r$ with $k<r\leq b$, since $u(a_q)>u(r)$ for all such $r$ and hence the reflection $t_{a_{m-1}a_m}$ would increase the length by more than $q$. Thus since in particular $u(r)>u(b)$ for all such $r$, we must have that $u(a_j)>u(b)$ for all $j\geq q$, as desired.
\end{proof}

\begin{lemma} \label{lemma:pieriknotkn1}
Suppose $u,w\in S_\infty$ and let $k>1$ be an integer. If $u\tom{k} w$ but $u\ntom{k-1} w$, then there is a unique $b>k$ such that $wt_{kb}\lup w$ and $u\tom{k-1} wt_{kb}$. Furthermore, in that case there does not exist a $a<k$ such that $wt_{ak}\lup w$ and $u\tom{k-1} wt_{ak}$.
\end{lemma}
\begin{proof}
Suppose $u\tom{k} w$ but $u\ntom{k-1} w$. Let $c_1,\ldots,c_n$ be the pairwise disjoint cycles in the decomposition of $u^{-1}w$, with the cycle length of $c_i$ being $p_i+1$. Some such cycle must contain $k$, as otherwise we would have $u\tom{k-1} w$ by Lemma \ref{lemma:piericycle}. Assume that $c_i=(a_{i,p_1},\ldots,a_{i,1},b_i)$ is the cycle that contains $k$ and set $b=b_i$. We claim that $u\tom{k-1} wt_{kb}$ and $wt_{kb}\lup w$. Note that $w=(wt_{kb})t_{kb}$ and hence 
$$(uc_1\cdots c_nt_{kb})t_{kb}=w$$
Let $j$ be the index such that $a_j=k$. Then we have
\begin{align*}
(a_{i,p_i},\ldots,a_{i,j+1},k,a_{i,j-1},\ldots,a_{i,1},b_i)t_{kb}&=(a_{i,j-1},\ldots,a_{i,1},b,a_{i,p_i},\ldots,a_{i,j+1},k)t_{kb}\\
&=(a_{i,j-1},\ldots,a_{i,1},b)(a_{i,p_i},\ldots,a_{i,j+1},k)
%&=t_{kb}(k,b)(a_{i,p_i},\ldots,a_{i,j+1},b)(a_{i,j-1},\ldots,a_{i,1},k)\\
%&=(a_{i,j-1},\ldots,a_{i,1},k)(a_{i,p_i},\ldots,a_{i,j+1},b)$$
\end{align*}
Note that $wt_{kb}(k)=w(b)=u(a_{i,p_i})$ and $wt_{kb}(b)=w(k)=u(a_{i,j-1})$. We need that
$$\ell(u,wt_{kb})=p_i'+p_i''-2(b-\desc_u(c_i'))-2(k-\desc_u(c_i''))+\ell(u,uc_1\cdots c_{i-1}c_{i+1}\cdots c_m)$$
We have that
$$\ell(u,w)=p_i-2(b-\desc_u(c_i))+\ell(u,uc_1\cdots c_{i-1}c_{i+1}\cdots c_m)$$
and
$$\ell(u,wt_{kb})\geq p_i'+p_i''-2(b-\desc_u(c_i'))-2(k-\desc_u(c_i''))+\ell(u,uc_1\cdots c_{i-1}c_{i+1}\cdots c_m)$$
If $c_i'$ has a descent, then $c_{i}''$ does not, and we are done. If $c_i''$ has a descent, then 
$$w(k)=u(a_{i,j-1})<u(b)<u(a_{i,p_i})=w(b)$$ 
by Lemma \ref{lemma:geq}, so 
$$\ell(w)<\ell(wt_{bk})\leq \ell(w)+2(b-k)-1$$
so
$$\ell(u,wt_{kb})=\ell(wt_{bk})-\ell(u)\leq p_i'+p_i''-2(b-\desc_u(c_i'))-2(b-\desc_u(c_i''))+2(b-k)+\ell(u,uc_1\cdots c_{i-1}c_{i+1}\cdots c_m)$$
Since we already have the opposite inequality, we are done if $c_{i}''$ has a descent, since $wt_{bk}\lup w$. If neither has a descent but $c_i$ does, then this descent must be at $k$. We have that $wt_{bk}(k)=u(a_{i,p_i})>u(b)$ by Lemma \ref{lemma:geq}, and $wt_{bk}(b)=u(a_{i,j-1})<u(b)$. Furthermore, $u(k)>u(b)>u(r)>u(a_{i,j-1})$ for all $r$ with $k<r<b$ since $k$ is a descent. Thus $wt_{bk}\lup w$.
%Suppose that there is a descent after $k$ in the cycle, say at %$%a_{i,q}$ for some $q> j$. This means that
%$$u(a_{i,q})>ut_{a_{i,1}b}\cdots ut_{a_{i,q-1}b}(r)>u(a_{i,q-1})$$
%for all $r$ with $a_{i,p_i}<r<b$. In particular, this means that
%$$u(a_{i,q})>ut_{a_{i,1}b}\cdots ut_{a_{i,q}b}(k)=w(k)=u(a_{i,j-1})>u(a_{i,p_i})=w(b)
%$$(a_{i_2},\ldots,a_{j+1},q')(a_{j-1},\ldots,a_{i_1},k)$$
%Note that $wt_{kb}(k)=u(b)$, $ut_{kb}(b)=u(k)$, and otherwise $ut_{kb}(a)=u(a)$ for all $a\notin \{k,b\}$. We have that
%$$\ell(u,w)=\sum_{i=1}(p_i-2(b_i-\desc_u(c_i)))$$
%If $u(k)>u(b)$, then
%$$\ell(ut_{kb},w)=-2(b-k)+p_1'+p_1''\sum_{i=2}(p_i-2(b_i-\desc_u(c_i)))$$
%$$ut_{a_{i,1}k}\cdots t_{a_{i,j-1}k}t_{a_{i,j+1}b}\cdots t_{a_{i,p_i}b}$$
%If $u(k)>u(b)$, then
%$$\ell(ut_{kb},w)=\ell(w)-\ell(u)+2(b-k)-1$$
%and otherwise
%$$\ell(ut_{kb},w)=\ell(w)-\ell(u)-2(k-\desc_u(c_1'))$$
%or
%$$\ell(ut_{kb},w)=\ell(w)-\ell(u)-2(b-\desc_u(c_1''))$$

%If $q<j$, then if $u(a_{i,q})>u(a_{i,q-1})$ we have that the length decreases by $2(k-a_{i,q})-1$.

%Suppose $q>j+1$. Then if $u(a_{i,q})>u(a_{i,q-1})$, we have that 
%$$u(a_{i,q})>ut_{a_{i,1}b}\cdots t_{a_{i,q-1}b}(r)>u(a_{i,q-1})$$
%for all 
$a_{i,q}<r<b$, so the length decreases by $2(b-a_{i,q})-1$ as desired. If $u(a_{i,j+1})>u(k)$

%If $u(k)>u(b)$, then $u(k)>u(a)>u(b)$ for all $k<a<b$.

 %Thus, if $l>j$ then any relation between $u(a_{i,l})$ and $u(a_{i,l+1})$ continues to hold, and $ut_{kb}(a_{i,j+1})=u(a_{i,j+1})$ has the same relationship to $ut_{kb}(b)=u(k)$. Further, $ut_{kb}(a_{i,1})>ut_{kb}(k)$ if and only if $u(a_{i,1})>u(b)$, so the relative length moves remain the same. If $l>j$ and $a_{i,l}\in \desc_u(c_i)$, then $a_{i,l}\in\desc_u(c_i')$. Necessarily $a_{i,l}<k$, and $u(a_{i,l})>u(a_{i,l-1})$. If $u(k)>u(b)$, then $\ell(ut_{kb})=\ell(u)-2(b-k)+1$. 
%, and the length decreases by some as yet undetermined $d_{i,l}$. If $1<l<j$ and $a_{i,l}\in\desc_u(c_i)$, then $a_{i,l}\in\desc_u(c_i'')$, and the length decreases by $d_{i,l}$.

To show that the $q'$ we have chosen is unique, let $r>k$ be any index other than $q'$ such that $\ell(ut_{kr})=\ell(u)+1$. Suppose first that $r$ is not in any of the cycles $c_i$. Recall that $c_1$ is the cycle that contains $k$. Then $t_{kr}c_1$ is a cycle in the disjoint cycle decomposition of $(ut_{kr})^{-1}w$ containing both $k$ and $r$, which are both larger than $k-1$. This violates the condition of Lemma \ref{lemma:piericycle} that in order for $ut_{kr}\tom{k-1} w$ we must have that each cycle contains exactly one element larger than $k-1$, hence we have $ut_{kr}\ntom{k-1} w$. If $r$ is in one of the cycles, since $r>k$ and $r\neq q'$ this cycle cannot be $c_1$. Say the cycle containing $r$ is $c_2=(a_{p_2}',\ldots,a_1',r)$. Then 
\begin{align*}
t_{kr}c_1c_2&=t_{kr}(r,a_{p_2}',\ldots,a_1')(a_{p_1},\ldots,a_1,q')\\
            &=(r,a_{p_2}',\ldots,a_1',k)(k,a_{j-1},\ldots,q',a_{p_1},\ldots,a_{j+1})\\
						&=(r,a_{p_2}',\ldots,a_1',k,a_{j-1},\ldots,q',a_{p_1},\ldots,a_{j+1})
\end{align*}
is a single cycle disjoint from all the others containing both $k$ and $r$,  which are both greater than $k-1$, hence we have $ut_{kr}\ntom{k-1} w$ by Lemma \ref{lemma:piericycle}. The uniqueness of $q'$ follows.

To prove the nonexistence of $q$, consider the same cycle decomposition $c_1,\ldots,c_n$. Let $q<k$ be such that $\ell(ut_{qk})=\ell(u)+1$.  If $c_1$ (the cycle that contains $k$) does not contain $q$, then $t_{qk}c_1$ is a cycle containing both $k$ and some integer larger than $k$. If some other cycle $c_2$ contains $q$, $t_{qk}c_1c_2$ is disjoint from all the other cycles and contains both $k$ and some integer larger than $k$, similarly to the calculation above. Thus if $c_1$ does not contain $q$ then the disjoint cycle decomposition of $(ut_{qk})^{-1}w$ contains the cycle $t_{qk}c_1$ or $t_{qk}c_1c_2$, which implies that $ut_{qk}\ntom{k-1} w$ by Lemma \ref{lemma:piericycle} since the cycle contains two elements larger than $k-1$. To handle the case where $c_1$ contains $q$, in that case $t_{qk}$ splits $c_1$ into two disjoint cycles as follows. Assume $c_1=(a_{p_1},\ldots,q')$ with $a_j=k$ and $a_{j'}=q$. We must have that $j'>j$ since $u(q)<u(k)$. Then 
\begin{align*}
t_{qk}c_1&=t_{qk}(a_{p_1},\ldots,q,\ldots,k,\ldots,q')\\
         &=t_{qk}(q,\ldots,k,\ldots,q',a_{p_1},\ldots,a_{j'+1})\\
				 &=t_{qk}(q,a_{j'-1},\ldots,k)(k,\ldots,q',a_{p_1},\ldots,a_{j'+1})\\
				 &=(q,a_{j'-1},\ldots,a_{j+1})(a_{p_1},\ldots,a_{j'+1},k,\ldots,q')
\end{align*}
The cycle $(a_{p_1},\ldots,k,\ldots,q')$ contains two elements larger than $k-1$, implying that $ut_{qk}\ntom{k-1} w$ by Lemma \ref{lemma:piericycle}, as desired.
\end{proof}

We organize the inductive step of Theorem \ref{theorem:pieri} around the position $k$, distinguishing whether or not $u(k)=w(k)$. The following lemma records the two combinatorial facts about the interaction of $\tom{k}$ and $\tom{k-1}$ at position $k$ that make this dichotomy exhaustive.

\begin{lemma} \label{lemma:leveldichotomy}
Suppose $u,w\in S_\infty$ and let $k>1$ be an integer.
\begin{enumerate}
\item[(i)] If $u(k)=w(k)$, then $u\tom{k}w$ if and only if $u\tom{k-1}w$.
\item[(ii)] If $u(k)\neq w(k)$ and $u\tom{k}w$, then $u\ntom{k-1}w$.
\end{enumerate}
\end{lemma}
\begin{proof}
Both statements follow from the cycle description of Lemma \ref{lemma:piericycle}, according to which $u\tom{m}w$ if and only if in the disjoint cycle decomposition of $u^{-1}w$ each cycle contains exactly one index exceeding $m$, subject to the length conditions. If $u(k)=w(k)$ then $k$ is a fixed point of $u^{-1}w$, hence lies in no nontrivial cycle; the constraint distinguishing level $k$ from level $k-1$ concerns only the indices exceeding $k-1$ that occur in the cycles, and $k$ is not among them. Thus the cycle conditions for level $k$ and level $k-1$ coincide, proving (i). For (ii), if $u(k)\neq w(k)$ then $k$ lies in a nontrivial cycle $c$ of $u^{-1}w$; if $u\tom{k}w$, the level-$k$ condition forces $k$ to be the unique index of $c$ exceeding $k-1$ that is admissible at level $k$, but then $c$ contains an index $b>k$ as its cycle top, so $c$ contains two indices $\geq k$, violating the level-$(k-1)$ condition. Hence $u\ntom{k-1}w$.
\end{proof}

The heart of the induction is the following cancellation, which replaces the informal ``the sums are empty'' step of the classical argument. When $u(k)=w(k)$ the diagonal and straight recursion terms already reproduce $c_k(p;u,w)$ through the single-variable recurrence for factorial elementary polynomials (Lemma \ref{lemma:elemrecursion}); the reflection sums and the quantum term must therefore cancel among themselves.

\begin{lemma}[Reflection--drop cancellation] \label{lemma:reflectioncancel}
Suppose $u,w\in S_\infty$, let $k>1$, and suppose $u(k)=w(k)$. Fix $1\leq p\leq k$. Then
$$\sum_{\substack{a<k,\ wt_{ak}\lup w}}\!\!\!\mathrm{sign}(a-k)\,\kappa_{ak}\,c_{k-1}(p-1;u,wt_{ak})\ +\ \sum_{\substack{b>k,\ wt_{kb}\lup w}}\!\!\!\mathrm{sign}(b-k)\,\kappa_{kb}\,c_{k-1}(p-1;u,wt_{kb})\ +\ q_{k-1}\,c_{k-2}(p-2;u,w)=0,$$
where $\kappa_{rs}=q_{rs}$ if $wt_{rs}\lup w$ is length-decreasing and $\kappa_{rs}=1$ if it is length-increasing. In fact the two reflection sums contribute at most one nonzero summand in total: it occurs precisely when $u\tom{k-2}w$ (equivalently, when the quantum term is nonzero), it is indexed by some $a<k$, and it equals $-\,q_{k-1}c_{k-2}(p-2;u,w)$.
\end{lemma}
\begin{proof}
Throughout we use that $u(k)=w(k)$, so that $k$ is a fixed point of $\pi:=u^{-1}w$.

\emph{Step 1: no reflection $t_{kb}$ with $b>k$ contributes.} Suppose $b>k$ and $wt_{kb}\lup w$ with $c_{k-1}(p-1;u,wt_{kb})\neq0$, so $u\tom{k-1}v$ for $v:=wt_{kb}$. Then $u^{-1}v=\pi t_{kb}$; since $k$ is a fixed point of $\pi$ and $b>k$, right multiplication by $t_{kb}$ adjoins $k$ to the $\pi$-cycle containing $b$ (or creates the $2$-cycle $(k\,b)$). Either way the resulting cycle of $u^{-1}v$ contains two indices exceeding $k-1$, namely $k$ and an index $\geq b>k$, contradicting the level-$(k-1)$ condition of Lemma \ref{lemma:piericycle}. Hence only reflections $t_{ak}$ with $a<k$ can contribute, and each such summand carries the sign $\mathrm{sign}(a-k)=-1$.

\emph{Step 2: the reflection sum has at most one nonzero term, occurring precisely when $u\tom{k-2}w$.} A reflection $t_{ak}$ with $a<k$ contributes precisely when $u\tom{k-1}v$ for $v:=wt_{ak}$; here $u^{-1}v=\pi t_{ak}$, so $t_{ak}$ adjoins the fixed point $k$ to the $\pi$-cycle containing $a$ (or creates the $2$-cycle $(a\,k)$ when $a$ is fixed by $\pi$). By Lemma \ref{lemma:piericycle}, $u\tom{k-1}v$ requires this enlarged cycle to contain exactly one index exceeding $k-1$; as it already contains $k$, it may contain no further index $\geq k$. Since $k$ is fixed by $\pi$, an index occurring in a nontrivial $\pi$-cycle and exceeding $k-1$ is in fact $\geq k+1$, unless it equals $k-1$. Thus the enlarged cycle is admissible only if the $\pi$-cycle receiving $k$ has all indices $\leq k-1$, i.e.\ its top equals $k-1$; the $\pi$-cycles being disjoint, at most one $a<k$ contributes. Such a configuration exists exactly when $u\tom{k-2}w$: if $u\tom{k-2}w$ and $k-1$ is not fixed by $\pi$, then the $\pi$-cycle through $k-1$ has top $k-1$ (any second index $\geq k-1$ would be $\geq k+1$, violating the level-$(k-2)$ condition), while if $k-1$ is fixed the choice $a=k-1$ creates the $2$-cycle $(k-1\,k)$; conversely a contributing $t_{ak}$ forces a $\pi$-cycle of top $k-1$ (or $k-1$ fixed with $a=k-1$), whence $u\tom{k-2}w$. This is exactly the condition under which the quantum term is nonzero.

\emph{Step 3: the contributing term cancels the quantum term.} Assume $u\tom{k-2}w$, let $t_{a^\ast k}$ ($a^\ast<k$) be the unique contributing reflection, and set $v:=wt_{a^\ast k}$. By Step 2 the cycle of $u^{-1}v$ containing $k$ has top $k$. Since $v$ agrees with $w$ off $\{a^\ast,k\}$ and $v(a^\ast)=w(k)=u(k)$, both $a^\ast$ and $k-1$ leave the agreement set in passing to level $k-1$, and comparison of the invariants of Definition~\ref{definition:pieri} gives
$$P_{k-1}(u,v)=P_{k-2}(u,w),\qquad n_{k-1}(u,v)=n_{k-2}(u,w)+1$$
(using $n_m=m-|P_m|$). Hence $c_{k-1}(p-1;u,v)$ and $c_{k-2}(p-2;u,w)$ carry the \emph{same} factorial elementary polynomial $E_{(p-2)-n_{k-2}(u,w),\,|P|}(y_P;z)$ with $P=P_{k-2}(u,w)$. Finally the $q$-weights match. Write $\kappa_{a^\ast k}=1$ if $t_{a^\ast k}$ is length-increasing and $\kappa_{a^\ast k}=q_{a^\ast k}$ if length-decreasing. As the $q$-weight is multiplicative over the cycles of $u^{-1}(\cdot)$ and only the cycle receiving $k$ is altered (from top $k-1$ to top $k$), the definition of the weights yields
$$\kappa_{a^\ast k}\,q_{k-1}(u,v)=q_{k-1}\,q_{k-2}(u,w).$$
Therefore the contributing summand is
$$\mathrm{sign}(a^\ast-k)\,\kappa_{a^\ast k}\,c_{k-1}(p-1;u,v)=-\,\kappa_{a^\ast k}\,q_{k-1}(u,v)\,E_P=-\,q_{k-1}\,q_{k-2}(u,w)\,E_P=-\,q_{k-1}c_{k-2}(p-2;u,w),$$
cancelling the quantum term. When $u\ntom{k-2}w$ the quantum term vanishes and, by Step 2, no reflection contributes; all three terms are $0$. In every case the displayed sum is $0$.
\end{proof}

\begin{remark} \label{remark:reflectionverify}
The content of Lemma \ref{lemma:reflectioncancel} is the identification of the quantum term $q_{k-1}c_{k-2}(p-2;u,w)$ with the unique level-$(k-1)$ neighbour $v$ satisfying $u\tom{k-1}v$, $u\ntom{k-2}v$, $P_{k-1}(u,v)=P_{k-2}(u,w)$ and $n_{k-1}(u,v)=n_{k-2}(u,w)+1$, obtained by adjoining the fixed point $k$ to the $\pi$-cycle of top $k-1$. The connecting reflection $t_{a^\ast k}$ may be length-increasing or length-decreasing; in either case its weight $\kappa_{a^\ast k}$ satisfies the balance $\kappa_{a^\ast k}q_{k-1}(u,v)=q_{k-1}q_{k-2}(u,w)$. Uniqueness of $v$ is a consequence of the cycle structure (there is a single $\pi$-cycle of top $k-1$); this intrinsic matching by the invariants $(P,n)$ is what the $u(k)\neq w(k)$ branch of Lemma \ref{lemma:vanishing} invokes.
\end{remark}

When $u(k)\neq w(k)$ the position $k$ lies in a nontrivial cycle of $\pi=u^{-1}w$, and the reflection sums of the vanishing identity below no longer telescope against the diagonal; instead they cancel among themselves and against the quantum term, one $q$-weight at a time. The following lemma isolates this cancellation. It is the moved-case counterpart of Lemma \ref{lemma:reflectioncancel}, and its non-quantum instances specialize at $M=1$ to the fixed-point-free reflection cancellation of the classical Pieri rule \cite{samuelmolev}.

\begin{lemma}[Moved reflection cancellation] \label{lemma:movedcancel}
Let $u,w\in S_\infty$, let $k>1$, and suppose $u\ntom{k}w$ with $u(k)\neq w(k)$. Let $\pi=u^{-1}w$, let $C$ be the $\pi$-cycle containing $k$ --- nontrivial, since $u(k)\neq w(k)$ --- and set $B=\max(C)$. Fix $1\leq p\leq k$ and a $q$-monomial $q^M\neq q_{k-1}(u,w)$, and let $\mathcal{G}_M$ be the set of Monk summands of \eqref{eq:up}--\eqref{eq:down} of $q$-weight $q^M$, augmented by the quantum term $q_{k-1}c_{k-2}(p-2;u,w)$ of \eqref{eq:rec} when $q^M=q_{k-1}q_{k-2}(u,w)$. Then the summands of $\mathcal{G}_M$ sum to $0$.
\end{lemma}
\begin{proof}
A neighbour $v=wt$, with $t\in\{t_{ak}:a<k\}\cup\{t_{kb}:b>k\}$, contributes to $\mathcal{G}_M$ only if $u\tom{k-1}v$; writing $\sigma=u^{-1}v=\pi t$, Lemma \ref{lemma:piericycle} requires every cycle of $\sigma$ to contain exactly one index exceeding $k-1$. Since $t$ moves $\pi$ only at its two positions, one of which is $k\in C$, the passage $\pi\rightsquigarrow\sigma$ either \emph{splits} $C$ (when the other position lies in $C$) or \emph{merges} $C$ with the component of the other position (when it does not); every other $\pi$-cycle is unchanged and so must already satisfy the level-$(k-1)$ condition.

\emph{Up-neighbours are unique and force $C$ to have exactly the indices $k,B$ above $k-1$.} Let $v=wt_{kb}$ with $b>k$. If $b\notin C$, the merge produces a cycle containing $k$, $b$ and $B\geq b$, hence two indices exceeding $k-1$: inadmissible. Thus $b\in C$, and $t_{kb}$ splits $C$ into a piece through $k$ and a piece through $b$; admissibility forces each to hold a single index above $k-1$, so $b$ is that index for its piece. As $B=\max(C)>k-1$, this yields $b=B$ and shows $C$ has exactly the two indices $k,B$ above $k-1$, with $B>k$. Hence a contributing up-neighbour, if present, is unique and equals $wt_{kB}$.

\emph{Down-neighbours.} Let $v=wt_{ak}$ with $a<k$. If $a\in C$ (a \emph{split}), admissibility assigns the indices of $C$ above $k-1$ one to each nontrivial piece: if $C$ has three or more such indices no split is admissible; if it has exactly two, $k$ and $B>k$, the split separates them (and then $wt_{kB}$ is the up-neighbour); if it has only $k$, so $B=k$, the split peels a fixed point off $C$. If $a\notin C$ (a \emph{merge}), then $C\cup(\text{component of }a)$ must retain a single index above $k-1$, namely $k$, forcing $B=k$ and the component of $a$ to have all indices $\leq k-1$. We record the resulting $q$-weights.

\emph{(1) The case $B=k$.} Either $a$ is adjoined as a fixed point --- then $v(a)=w(k)=u(a)$ reinstates below $k$ the agreement lost at $k$, the cycle tops are unchanged, and the summand carries weight $q_{k-1}(u,w)$ --- or the component absorbed into $C$ is a nontrivial cycle, whose unique index above $k-2$ is then $k-1$ (being $\leq k-1$ and $>k-2$); the merged cycle acquires the extra factor $q_{k-1}$, and a direct comparison of the invariants of Definition \ref{definition:pieri} gives
$$P_{k-1}(u,v)=P_{k-2}(u,w),\qquad n_{k-1}(u,v)=n_{k-2}(u,w)+1,\qquad \kappa_{ak}q_{k-1}(u,v)=q_{k-1}q_{k-2}(u,w),$$
the very identities of Lemma \ref{lemma:reflectioncancel}, with the $\pi$-cycle of top $k-1$ now attached to $C$ in place of the fixed point $k$; the computation is otherwise verbatim, and by uniqueness of the $\pi$-cycle of top $k-1$ there is at most one such $v$.

\emph{(2) The case $B>k$.} Here only the split with $C=\{k,B\}$ above $k-1$ occurs, and it separates $k$ from $B$.

\emph{Cancellation.} Suppose first $q^M=q_{k-1}q_{k-2}(u,w)$, so the quantum term lies in $\mathcal{G}_M$ and $u\tom{k-2}w$; then every $\pi$-cycle has exactly one index above $k-2$, in particular $B=k$, and by the up-paragraph $\mathcal{G}_M$ contains no up-neighbour. By (1) the summands of weight $q^M$ reduce to the single merge down-neighbour $v$ with $P_{k-1}(u,v)=P_{k-2}(u,w)$; then $c_{k-1}(p-1;u,v)$ and $c_{k-2}(p-2;u,w)$ carry the same factorial elementary polynomial $E_{(p-2)-n_{k-2}(u,w),\,(k-2)-n_{k-2}(u,w)}(y_{P_{k-2}(u,w)};z)$, and the down summand, of sign $\mathrm{sign}(a-k)=-1$, equals $-q_{k-1}c_{k-2}(p-2;u,w)$ and cancels the quantum term. (If no such $v$ exists the quantum term already vanishes.)

Now suppose $q^M\neq q_{k-1}(u,w),\,q_{k-1}q_{k-2}(u,w)$. By (1) every merge or peel down carries weight $q_{k-1}(u,w)$ or $q_{k-1}q_{k-2}(u,w)$, hence lies outside $\mathcal{G}_M$; so every summand of $\mathcal{G}_M$ is of split type with $C=\{k,B\}$, $B>k$ (whence also $u\ntom{k-1}w$). For each such summand $P_{k-1}(u,v)=P_{k-1}(u,w)$: the up $wt_{kB}$ agrees with $w$ below $k$, while a split down $wt_{ak}$ agrees with $w$ below $k$ except at $a\in C$, where it takes the value $w(k)$; but $a$ is not fixed by $\pi$, and $u(a)\neq w(k)$ (otherwise, by (1), the summand would have weight $q_{k-1}(u,w)$), so no agreement below $k$ is created or destroyed. Consequently all summands of $\mathcal{G}_M$ share the factorial elementary factor $E_0=E_{(p-1)-n,\,(k-1)-n}(y_{P};z)$, where $P=P_{k-1}(u,w)$ and $n=n_{k-1}(u,w)$. Finally, when $C=\{k,B\}$ there are exactly two admissible neighbours through position $k$: the up $wt_{kB}$, and --- since $C$ has a single descent (Lemma \ref{lemma:onedescent}) and, as in the proof of Lemma \ref{lemma:pieriknotkn1}, exactly one covering split of $C$ separates $k$ from $B$ --- a single down $wt_{a^\ast k}$. The two reflections differ by the $3$-cycle $(a^\ast,B,k)$; each resolves the cycle $C$ into two cycles of tops $k$ and $B$ while leaving the other cycles untouched, so by the multiplicativity of the $q$-weight (Definition \ref{definition:pieri}) and Lemma \ref{lemma:geq}, which fixes $\desc_u(C)$ and hence the drop factors, both $\kappa_{kB}q_{k-1}(u,wt_{kB})$ and $\kappa_{a^\ast k}q_{k-1}(u,wt_{a^\ast k})$ equal the weight $q^M$ of $\mathcal{G}_M$. With opposite signs $\mathrm{sign}(B-k)=+1$ and $\mathrm{sign}(a^\ast-k)=-1$ and common coefficient $E_0$, the up and the down cancel. Hence $\mathcal{G}_M=0$.
\end{proof}

The remaining case of the induction is when $u\ntom{k}w$, where the target coefficient is $0$.

\begin{lemma}[Vanishing] \label{lemma:vanishing}
Suppose $u,w\in S_\infty$, let $k>1$, and suppose $u\ntom{k}w$. Then for all $1\leq p\leq k$,
\begin{align*}
(y_{w(k)}-z_{k-p+1})\,c_{k-1}(p-1;u,w)
&+\sum_{\substack{b>k,\ wt_{kb}\lup w}}\kappa_{kb}\,c_{k-1}(p-1;u,wt_{kb})
-\sum_{\substack{a<k,\ wt_{ak}\lup w}}\kappa_{ak}\,c_{k-1}(p-1;u,wt_{ak})\\
&+c_{k-1}(p;u,w)+q_{k-1}c_{k-2}(p-2;u,w)=0,
\end{align*}
where $\kappa_{rs}=q_{rs}$ if $wt_{rs}\lup w$ is length-decreasing and $\kappa_{rs}=1$ if it is length-increasing.
\end{lemma}
\begin{proof}
Suppose first that $u(k)=w(k)$. Then by Lemma \ref{lemma:leveldichotomy}(i) also $u\ntom{k-1}w$, so both $c_{k-1}(p-1;u,w)$ and $c_{k-1}(p;u,w)$ vanish; hence the diagonal term and the straight recursion term are $0$, and the remaining reflection sums and the quantum term cancel by Lemma \ref{lemma:reflectioncancel}. This settles the case $u(k)=w(k)$ completely.

Now suppose $u(k)\neq w(k)$. Since $u\ntom{k}w$, Lemma \ref{lemma:piericycle} shows the cycle of $u^{-1}w$ containing $k$ has a top $b>k$, so $c_k(p;u,w)=0$; writing $F$ for the right-hand side \eqref{eq:diag}--\eqref{eq:rec}, we must show $F=0$.

We grade by the cohomological degree $\deg q_i=2$. Under this grading \eqref{eq:pieriexpand} is homogeneous, and the $q_i$ are algebraically independent over $\mathbb{Z}[y,z]$. By the definition of $c_m$, each summand of $F$ is a monomial in the $q_i$ times a factorial elementary polynomial in $\mathbb{Z}[y,z]$; collecting by $q$-monomial we write $F=\sum_M q^M F_M$ with $F_M\in\mathbb{Z}[y,z]$, and
\[
F=0 \iff F_M=0 \text{ for every } q\text{-monomial } M .
\]
Fix $M$ and let $\mathcal{G}_M$ be the collection of summands of $F$ whose $q$-weight equals $q^M$; here the diagonal \eqref{eq:diag} and the straight term $c_{k-1}(p;u,w)$ carry weight $q_{k-1}(u,w)$, the quantum term carries $q_{k-1}q_{k-2}(u,w)$, and the Monk summand of \eqref{eq:up}--\eqref{eq:down} reaching $w'=wt_{pp,k}$ carries $\kappa\,q_{k-1}(u,w')$ with $\kappa\in\{1,q_{a k},q_{k b}\}$ its drop-weight. We show $F_M=0$ in two exhaustive cases.

\emph{Case (i): $\mathcal{G}_M$ contains the diagonal or the straight term.} These carry weight $q_{k-1}(u,w)$, so $q^M=q_{k-1}(u,w)$; in particular $u\tom{k-1}w$. Write $P=P_{k-1}(u,w)$, $n=n_{k-1}(u,w)$, $m=k-1-n=|P|$ and $d=p-1-n$. Set $a_0=u^{-1}(w(k))$; since $u(k)\neq w(k)$ and $w(k)$ is a value of $u$ occurring at a position $<k$ (this holds whenever the diagonal is nonzero), we have $a_0<k$. The Monk summand indexed by $t_{a_0 k}$ then satisfies $u\tom{k-1}wt_{a_0 k}$ with
$$P_{k-1}(u,wt_{a_0k})=P\cup\{w(k)\},\qquad n_{k-1}(u,wt_{a_0k})=n-1,$$
carries weight $q_{k-1}(u,w)$, and is the \emph{only} Monk summand in $\mathcal{G}_M$. Dividing $\mathcal{G}_M$ by the common factor $q_{k-1}(u,w)$, its three contributions read
$$(y_{w(k)}-z_{k-p+1})\,E_{d,m}(y_P;z)+E_{d+1,m}(y_P;z)-E_{d+1,m+1}(y_{P\cup\{w(k)\}};z),$$
which vanishes by the single-variable recurrence of Lemma \ref{lemma:elemrecursion} adjoining the variable $y_{w(k)}$ (the $z$-index matches because $m-d=k-p$). Hence $F_M=0$.

\emph{Case (ii): $\mathcal{G}_M$ contains neither the diagonal nor the straight term.} Then $q^M\neq q_{k-1}(u,w)$, and $\mathcal{G}_M$ consists of Monk summands of \eqref{eq:up}--\eqref{eq:down} together with the quantum term when $q^M=q_{k-1}q_{k-2}(u,w)$. By Lemma \ref{lemma:movedcancel} these summands sum to $0$, so $F_M=0$.

As the two cases are exhaustive, every $F_M=0$, so $F=0$.
\end{proof}

We now record the recurrence satisfied by the quantum factorial elementary polynomials, which drives the induction.

\begin{lemma} \label{lemma:elemrecursion}
For all integers $p$ and $k$ we have
$$E_{p,k}^q(x;z)=(x_k-z_{k-p+1})E_{p-1,k-1}^q(x;z)+E_{p,k-1}^q(x;z)+q_{k-1}E_{p-2,k-2}^q(x;z),$$
with the conventions $E_{0,k}^q=1$ for $k\geq 0$ and $E_{p,k}^q=0$ if $p<0$ or $p>k$. Moreover, the ordinary factorial elementary polynomials satisfy the single-variable recurrence
$$E_{a,m}(y_1,\ldots,y_m;z)=E_{a,m-1}(y_1,\ldots,y_{m-1};z)+(y_m-z_{m-a+1})E_{a-1,m-1}(y_1,\ldots,y_{m-1};z).$$
\end{lemma}

The first recurrence is the quantization of the classical relation $E_{p,k}=(x_k-z_{k-p+1})E_{p-1,k-1}+E_{p,k-1}$; the extra term $q_{k-1}E_{p-2,k-2}^q$ is the only feature absent from the non-quantum setting. The classical Pieri formula, obtained by discarding all quantum terms, is proved in \cite{samuelmolev} by induction on $k$ using the first recurrence in the case $q=0$; we adapt that argument, the essential new point being the bookkeeping of the powers of $q$.

\begin{proof}[Proof of Theorem \ref{theorem:pieri}]
For a permutation $w$, write $[\sch_w^q]F$ for the coefficient of $\sch_w^q(x;y)$ in an expansion $F$ of a quantum double Schubert polynomial in the basis $\{\sch_w^q(x;y)\}$. With this notation the theorem is the assertion that, for all $u$ and all $0\leq p\leq k$,
$$[\sch_w^q]\big(\sch_u^q(x;y)E_{p,k}^q(x;z)\big)=\begin{cases}q_k(u,w)\,E_{p-n_k(u,w),\,k-n_k(u,w)}(y_{P_k(u,w)};z) & \text{if }u\tom{k} w,\\[2pt] 0 & \text{otherwise.}\end{cases}$$
For brevity, when $u\tom{m} w$ we set
$$c_m(a;u,w)=q_m(u,w)\,E_{a-n_m(u,w),\,m-n_m(u,w)}(y_{P_m(u,w)};z),$$
and $c_m(a;u,w)=0$ if $u\ntom{m}w$; note that $c_m(a;u,w)=0$ whenever $a<n_m(u,w)$ or $a>m$, since then the factorial elementary factor vanishes. The claim is that $[\sch_w^q](\sch_u^qE_{p,k}^q)=c_k(p;u,w)$.

We argue by strong induction on $k$, the statement for level $k$ being understood to include all $p$ simultaneously.

\emph{Base case $k=1$.} Only $p\in\{0,1\}$ give a nonzero $E_{p,1}^q$. For $p=0$ we have $E_{0,1}^q=1$ and the claim is trivial, since $u\tom{1}w$ with $q_1=1$, $n_1=0$ only for $w=u$. For $p=1$ we have $E_{1,1}^q(x;z)=x_1-z_1$, and Lemma \ref{lemma:singlevariable} with $i=j=1$ gives
$$(x_1-z_1)\sch_u^q(x;y)=(y_{u(1)}-z_1)\sch_u^q(x;y)+\sum_{\substack{b>1\\\ell(ut_{1b})=\ell(u)+1}}\sch_{ut_{1b}}^q(x;y)+\sum_{\substack{b>1\\\ell(ut_{1b})=\ell(u)-2(b-1)+1}}q_{1b}\sch_{ut_{1b}}^q(x;y).$$
The permutations $w$ with $u\tom{1}w$ are exactly $w=u$ and $w=ut_{1b}$ for $b>1$ with $u\lup ut_{1b}$, since Definition \ref{definition:pieri} permits at most one reflection when $k=1$. For $w=u$ we have $q_1=1$, $n_1=0$, $P_1(u,u)=\{u(1)\}$, and $E_{1,1}(y_{u(1)};z)=y_{u(1)}-z_1$, matching the diagonal term. For $w=ut_{1b}$ we have $n_1=1$, $P_1=\emptyset$, and $E_{0,0}(\,;z)=1$; if the step is length-increasing then $q_1(u,w)=1$, matching the second sum, and if it is length-decreasing then $q_1(u,w)=q_{1b}$, matching the third sum. All other coefficients vanish on both sides. This proves the base case.

\emph{Inductive step.} Fix $k>1$ and $1\leq p\leq k$, and assume the claim at all levels $m<k$. By Lemma \ref{lemma:elemrecursion},
\begin{equation}\label{eq:pieriexpand}
\sch_u^qE_{p,k}^q=(x_k-z_{k-p+1})\big(\sch_u^qE_{p-1,k-1}^q\big)+\sch_u^qE_{p,k-1}^q+q_{k-1}\,\sch_u^qE_{p-2,k-2}^q.
\end{equation}
By the inductive hypothesis, $\sch_u^qE_{p-1,k-1}^q=\sum_v c_{k-1}(p-1;u,v)\,\sch_v^q$. Applying Lemma \ref{lemma:singlevariable} with $i=k$, $j=k-p+1$ to each $\sch_v^q$ and extracting the coefficient of $\sch_w^q$ from \eqref{eq:pieriexpand}, we obtain
\begin{align}
[\sch_w^q]\big(\sch_u^qE_{p,k}^q\big)
&=(y_{w(k)}-z_{k-p+1})\,c_{k-1}(p-1;u,w) \label{eq:diag}\\
&\quad+\sum_{\substack{b>k,\ w t_{kb}\lup w\\ \ell(w)=\ell(wt_{kb})+1}} c_{k-1}(p-1;u,wt_{kb})
+\sum_{\substack{b>k,\ w t_{kb}\lup w\\ \ell(w)=\ell(wt_{kb})-2(b-k)+1}} q_{kb}\,c_{k-1}(p-1;u,wt_{kb}) \label{eq:up}\\
&\quad-\sum_{\substack{a<k,\ w t_{ak}\lup w\\ \ell(w)=\ell(wt_{ak})+1}} c_{k-1}(p-1;u,wt_{ak})
-\sum_{\substack{a<k,\ w t_{ak}\lup w\\ \ell(w)=\ell(wt_{ak})-2(k-a)+1}} q_{ak}\,c_{k-1}(p-1;u,wt_{ak}) \label{eq:down}\\
&\quad+c_{k-1}(p;u,w)+q_{k-1}\,c_{k-2}(p-2;u,w), \label{eq:rec}
\end{align}
where we used that $\mathrm{sign}(b-k)=1$ for $b>k$ and $\mathrm{sign}(a-k)=-1$ for $a<k$, and that a summand in \eqref{eq:up}--\eqref{eq:down} is present precisely when $u\tom{k-1}wt_{kb}$, respectively $u\tom{k-1}wt_{ak}$ (otherwise $c_{k-1}=0$). We must show this equals $c_k(p;u,w)$. We distinguish three cases according to whether $u(k)=w(k)$ and, when $u(k)\neq w(k)$, whether $u\tom{k}w$; by Lemma \ref{lemma:leveldichotomy} these cases are exhaustive and, in particular, the combination $u(k)\neq w(k)$ with both $u\tom{k}w$ and $u\tom{k-1}w$ never occurs.

\emph{Case 1: $u(k)=w(k)$.} By Lemma \ref{lemma:leveldichotomy}(i) we have $u\tom{k}w$ if and only if $u\tom{k-1}w$; suppose first that both hold. Since $u(k)=w(k)$,
$$n_k(u,w)=n_{k-1}(u,w),\qquad P_k(u,w)=P_{k-1}(u,w)\cup\{u(k)\},\qquad q_k(u,w)=q_{k-1}(u,w),$$
the last equality holding because no reflection with first index $k$ occurs. Writing $m=k-n_k(u,w)=|P_k(u,w)|$ and $a=p-n_k(u,w)$, the diagonal contribution \eqref{eq:diag} and the straight recursion term $c_{k-1}(p;u,w)$ of \eqref{eq:rec} give
$$(y_{u(k)}-z_{k-p+1})\,c_{k-1}(p-1;u,w)+c_{k-1}(p;u,w)=q_{k-1}(u,w)\Big[(y_{u(k)}-z_{m-a+1})E_{a-1,m-1}(y_{P_{k-1}};z)+E_{a,m-1}(y_{P_{k-1}};z)\Big],$$
using $k-p+1=m-a+1$. By the single-variable recurrence of Lemma \ref{lemma:elemrecursion} with the new variable $y_{u(k)}$, the bracket equals $E_{a,m}(y_{P_k};z)$, so these two contributions sum to $q_k(u,w)E_{a,m}(y_{P_k};z)=c_k(p;u,w)$. The remaining contributions are the reflection sums \eqref{eq:up}, \eqref{eq:down} and the quantum term $q_{k-1}c_{k-2}(p-2;u,w)$ of \eqref{eq:rec}, which sum to $0$ by Lemma \ref{lemma:reflectioncancel}. Hence the total is $c_k(p;u,w)$, as required. If instead $u(k)=w(k)$ but $u\ntom{k}w$ (equivalently $u\ntom{k-1}w$), then $c_k(p;u,w)=0$ and the diagonal and straight terms both vanish, while the reflection sums and quantum term again cancel by Lemma \ref{lemma:reflectioncancel}; the total is $0=c_k(p;u,w)$.

\emph{Case 2: $u\tom{k}w$ but $u\ntom{k-1}w$.} By Lemma \ref{lemma:pieriknotkn1} there is a unique $b>k$ with $wt_{kb}\lup w$ and $u\tom{k-1}wt_{kb}$; set $v=wt_{kb}$. This is exactly the unique nonvanishing summand of \eqref{eq:up}, and by Lemma \ref{lemma:pieriknotkn1} the sums \eqref{eq:down} vanish, as does the diagonal \eqref{eq:diag} and the terms \eqref{eq:rec} (since $u\ntom{k-1}w$ and $u\ntom{k-2}w$). Moreover $P_k(u,w)=P_{k-1}(u,v)$ and $n_k(u,w)=n_{k-1}(u,v)+1$, so
$$E_{p-n_k(u,w),\,k-n_k(u,w)}(y_{P_k};z)=E_{(p-1)-n_{k-1}(u,v),\,(k-1)-n_{k-1}(u,v)}(y_{P_{k-1}(u,v)};z).$$
If $\ell(w)=\ell(v)+1$ then the reflection $t_{kb}$ is length-increasing, so by Definition \ref{definition:pieri} $q_k(u,w)=q_{k-1}(u,v)$ and the contribution is $c_{k-1}(p-1;u,v)=c_k(p;u,w)$. If instead $\ell(w)=\ell(v)-2(b-k)+1$ then $t_{kb}$ is length-decreasing, Definition \ref{definition:pieri} gives $q_k(u,w)=q_{k-1}(u,v)\,q_{kb}$, and the contribution is $q_{kb}\,c_{k-1}(p-1;u,v)=c_k(p;u,w)$. In both cases the coefficient equals $c_k(p;u,w)$, as required. This is the step at which the quantum drop terms of Lemma \ref{lemma:singlevariable} and the recursive definition of $q_k(u,w)$ interlock.

\emph{Case 3: $u\ntom{k}w$ (with $u(k)\neq w(k)$).} Here $c_k(p;u,w)=0$, so we must show the coefficient \eqref{eq:diag}--\eqref{eq:rec} vanishes. This is exactly the assertion of Lemma \ref{lemma:vanishing}. (The subcase $u(k)=w(k)$ with $u\ntom{k}w$ was already treated at the end of Case 1.)

In all cases $[\sch_w^q](\sch_u^qE_{p,k}^q)=c_k(p;u,w)$, completing the induction and the proof.
\end{proof}

\bibliographystyle{acm}
\bibliography{formschub}
\end{document}

